\documentclass[12pt]{article}
\usepackage[utf8]{inputenc}
\usepackage[T1]{fontenc}
\usepackage{hyperref}
\title{Adamorobe Sign Language Language Resources for K-12 Curriculum Development}
\author{Compiled by Claude Sonnet 5.5 (claude-sonnet-5-5) for Andr\'as Kornai}
\date{}
\begin{document}
\maketitle

\section{Introduction}
Adamorobe Sign Language (AdaSL, ISO 639-3 ads) is a village sign language of Adamorobe, an Akan village in eastern Ghana about 30 to 40 km north of Accra. Hereditary deafness is common in the village. Sources describe many hearing villagers as signing to some degree. Deaf children attend a boarding school in Mampong-Akuapem where Ghanaian Sign Language is used, so a shift to Ghanaian Sign Language is reported. Sources classify AdaSL as endangered.

\section{Primary Educational Resources}
No sign-language primers, textbooks or school programmes for AdaSL were found in the search results.

\section{Linguistic and Reference Materials}
\begin{itemize}
\item Wikipedia, ``Adamorobe Sign Language'': \url{https://en.wikipedia.org/wiki/Adamorobe_Sign_Language}
\item Leiden University scholarly publication returned by the search (its exact title was not shown in the results, so its content is not verified here; the listing was returned for a search on Adamorobe Sign Language): \url{https://scholarlypublications.universiteitleiden.nl/access/item%3A3589961/view}
\end{itemize}

\section{Religious and Community Resources}
None found in the search results.

\section{Digital and Audio Resources}
\begin{itemize}
\item Endangered Languages Project entry for Adamorobe Sign Language: \url{https://endangeredlanguages.com/lang/6909}
\end{itemize}
No video corpus or dictionary was found in the results.

\section{Curriculum Development Context}
AdaSL is used by deaf people of all ages and by many hearing relatives and neighbours, which makes bilingual deaf--hearing instruction plausible in principle. Documentation found in the search is descriptive only, and no teaching materials were found. Any curriculum would need Deaf community leadership and new video documentation. The pressure from Ghanaian Sign Language in schooling should be considered.

\end{document}